% mode: ini
% WEB 1461 and 1462: \errmessage with \errhelp set: use_err_help, so error calls give_err_help, in the mode where help goes on
% the transcript.  An empty \errhelp means null, so the usual help comes back: its first use counts for long_help_seen even
% with \errhelp in force before it (not set in this file: the first \errmessage uses \errhelp, and only the first without it
% sees the long help).  \newlinechar inside \errhelp, which error sets for the help.  A second \errmessage with \errhelp is not
% "another".  An error of another kind right after: use_err_help has gone back to false.  \errhelp in a group.
\input prelude
\errorcontextlines=0
\errhelp={help one}
\errmessage{with help}
\errmessage{with help again}
\errhelp={}
\errmessage{long help now}
\errmessage{the short one}
\errhelp={help two}
\errmessage{help again}
\errhelp={}
\errmessage{the short one still}
\errhelp={line 1^^Jline 2}
\errmessage{no newlinechar}
\newlinechar=`^^J
\errmessage{with newlinechar}
\newlinechar=-1
\errhelp={\relax # \tenrm ~ {a} $}
\errmessage{tokens in the help}
\undefinedone
\errhelp={a long help text, longer than one line of the terminal, which is seventy nine characters wide}
\errmessage{long}
\begingroup\errhelp={local}\errmessage{in a group}\endgroup
\errmessage{after the group}
\global\errhelp={}
\errmessage{void}
\toks0={from a register}\errhelp=\toks0 \errmessage{from a register}
\errhelp={}
\scrollmode
\errhelp={in scroll mode}
\errmessage{scroll}
\errhelp={}
\errmessage{scroll, without}
\batchmode
\errhelp={in batch mode}
\errmessage{batch}
\nonstopmode
\errhelp={}
\end
